% -----------------------------------------------------------------------------
% chapters/chap-algorithm.tex —— 算法流程案例：线性规划的单纯形法
%
% 本文件演示“把一章单独放进 chapters/ 目录、由主文件 \input 调用”的写法。
% 文件本身以 \chapter 开头，因此被调用后即成为书中的一章。
%
% 注意：本例只需要 booktabs 表格环境，不需要额外的算法宏包；
%       若要使用 algorithm2e、algorithmicx 等算法环境，请在主文件导言区加载，
%       而不是在类文件中加载。
% -----------------------------------------------------------------------------
\chapter{算法流程案例：线性规划的单纯形法}

本章以线性规划为例，说明如何把一个算法的执行流程写进书稿。算法流程的书写没有强制格式，
本模板提供的基本工具已经够用：\textbf{有序列表}用于描述步骤，\textbf{表格}用于展示中间结果，
\textbf{定理类环境}用于给出结论与证明，\lstinline{note} 环境用于补充说明。

\section{问题描述}

考虑如下线性规划问题：

\begin{equation}
  \begin{aligned}
    \max \quad & z = 3x_1 + 5x_2 \\
    \mathrm{s.t.}\quad & x_1 \leq 4, \\
                        & 2x_1 + 3x_2 \leq 18, \\
                        & x_1, x_2 \geq 0.
  \end{aligned}
  \label{eq:lp}
\end{equation}

目标是求使 $z$ 最大的可行解。可行域是一个多边形，最优解必然出现在某个顶点上，
但当约束较多时逐一枚举顶点并不可行，这正是单纯形法的用武之地。

\section{算法流程}

单纯形法的完整流程如下，每一步都对应表格的一次更新：

\begin{enumerate}
  \item \textbf{化为标准型}：引入松弛变量 $x_3, x_4$，把不等式约束改写为等式
    $x_1 + x_3 = 4$、$2x_1 + 3x_2 + x_4 = 18$，并把目标函数写为
    $z - 3x_1 - 5x_2 = 0$。以 $x_3, x_4$ 为初始基，写出初始单纯形表。
  \item \textbf{计算检验数}：观察目标函数行（$z$ 行）中各非基变量的系数。
    只要存在负系数，就说明当前解不是最优解。
  \item \textbf{选择入基变量}：取检验数最小（最负）的非基变量所在列。
    本例中 $x_2$ 的检验数为 $-5$，故 $x_2$ 入基。
  \item \textbf{计算比值并选择基变量}：对入基列中系数为正的各行计算比值
    $b_i / a_{ij}$，取最小者；该行对应的基变量离基。这一步保证主元为正，
    从而维持可行性。若某行该列系数不大于零，则不参与比值。
  \item \textbf{行变换（高斯消元）}：以主元为支点，把主元化为 $1$，再把同列其余元素
    化为 $0$，得到新的单纯形表。
  \item \textbf{判断是否最优}：若 $z$ 行不再有负系数，则当前基本可行解即为最优解；
    否则回到第 2 步继续迭代。
  \item \textbf{输出结果}：读出基变量取值与非基变量取值，代回目标函数得到最优值。
\end{enumerate}

\begin{note}
单纯形法的每一步迭代都保持可行性（比值检验的作用），但不能保证有限步终止——
退化情形下可能出现循环。因此实际应用中通常配合 \textbf{ Bland 规则} 等防循环策略。
\end{note}

\section{初始单纯形表与迭代}

初始时 $x_3, x_4$ 为基变量，非基变量为 $x_1, x_2$：

\begin{table}[htbp]
  \centering
  \caption{初始单纯形表}
  \label{tab:simplex-0}
  \begin{tabular}{c|ccccc}
    \toprule
    & $x_1$ & $x_2$ & $x_3$ & $x_4$ & $b$ \\
    \midrule
    $x_3$ & 1 & 0 & 1 & 0 & 4 \\
    $x_4$ & 2 & 3 & 0 & 1 & 18 \\
    \midrule
    $z$  & $-3$ & $-5$ & 0 & 0 & 0 \\
    \bottomrule
  \end{tabular}
\end{table}

$z$ 行中 $-5$ 是最小的负数，故 $x_2$ 入基。比值检验给出
$\frac{4}{0} = +\infty$（不参与）与 $\frac{18}{3} = 6$，故 $x_4$ 离基，
主元为第二行第三列的 $3$。行变换后得到：

\begin{table}[htbp]
  \centering
  \caption{迭代一次后的单纯形表}
  \label{tab:simplex-1}
  \begin{tabular}{c|ccccc}
    \toprule
    & $x_1$ & $x_2$ & $x_3$ & $x_4$ & $b$ \\
    \midrule
    $x_3$ & 1 & 0 & 1 & 0 & 4 \\
    $x_2$ & $\frac{2}{3}$ & 1 & 0 & $\frac{1}{3}$ & 6 \\
    \midrule
    $z$  & $\frac{1}{3}$ & 0 & 0 & $\frac{5}{3}$ & 30 \\
    \bottomrule
  \end{tabular}
\end{table}

此时 $z$ 行已无非负项以外的负数，迭代结束。

\section{结论}

\begin{conclusion}
问题 \eqref{eq:lp} 的最优解为 $x_1 = 0$、$x_2 = 6$，最优值为 $z = 30$。
\end{conclusion}

\begin{exercise}
用单纯形法求解下述线性规划，并写出所用的单纯形表：
\begin{equation}
  \max\; z = 2x_1 + 3x_2 \quad
  \mathrm{s.t.}\; x_1 + x_2 \leq 5,\; x_1 + 2x_2 \leq 8,\; x_1, x_2 \geq 0.
\end{equation}
\end{exercise}
